\section{Established computational results}

\subsection{Identity-safe local cascade}

At fixed $(a,c)=(0.245,5.1)$, the true period-5 family undergoes seven resolved
supercritical flips.  The first six parent-event values are
\begin{equation}
\begin{split}
b_1&=0.1834675907716,\quad b_2=0.1805372082024,\\
b_3&=0.1798912237616,\quad b_4=0.1797506213663,\\
b_5&=0.1797203688505,\quad b_6=0.1797138833005,
\end{split}
\end{equation}
and the resolved period-320 to period-640 event is
$b_7=0.17971249399393$.  The associated finite spacing ratios are
$4.5363$, $4.5944$, $4.6476$, $4.6646$, and $4.6682$.  Multiple shooting,
block-Floquet stability, and whole-orbit identity establish a stable period-640
child.  These are finite numerical observations; proximity of the ratios to a
classical constant is not presented as a universality proof.

Three event locations were predicted from earlier rungs before their dedicated
search.  The frozen period-80 to period-160 prediction missed the resolved
event by $1.398\times10^{-7}$.  The frozen period-320 to period-640 prediction
missed by $3.00\times10^{-10}$.  The period-640 to period-1280 prediction at
$b=0.1797121964470$ prospectively brackets a real $-1$ crossing between
$b=0.17971219$ and $b=0.17971220$; the bracket midpoint differs from the
prediction by $1.447\times10^{-9}$.  This third test has not yet passed event
refinement or period-1280 child qualification.

Subsequent baseline and tighter-solver scalar refinements retain a signed
interval only $3.22\times10^{-15}$ wide, but the closest precision-audited
point has multiplier residual $1.70\times10^{-8}$ and therefore misses the
preregistered $10^{-8}$ equality gate.  The block-cyclic multiplier and direct
monodromy products agree near $10^{-13}$, so the interval is retained as strong
event evidence but is not promoted to a corrected eighth event.  A coupled
orbit--parameter--eigenvector solve with an exact second-variational Jacobian
recovers the already resolved seventh event within $8.94\times10^{-13}$, with
orbit and tangent residuals $9.21\times10^{-13}$ and $1.01\times10^{-11}$.
This validates the event formulation on a known case; it must still pass at the
period-640 parent before branch switching.

\subsection{Flow bifurcation versus section geometry}

After the period-5 to period-10 flip, the stable period-10 child changes its
crossing count on the historical return section near
$b=0.181750232321$.  Continuous tangency refinement identifies this as a
section-boundary grazing.  It is therefore not another bifurcation of the flow
orbit.  This result is directly relevant to branch-count studies: an apparent
change in a scalar return plot can originate in the measurement construction,
so a topology-transition classifier must include section-robustness tests.

\subsection{Long transient capture}

Cases that initially appeared to show coexistence of a chaotic attractor and a
stable periodic attractor were followed for longer horizons and with multiple
initial conditions.  The trajectories eventually entered the stable periodic
windows.  Persistent multistability is rejected for those tested cases; the
evidence instead supports long chaotic-transient capture, consistent with a
nonattracting chaotic-saddle mechanism.  Computing the saddle itself remains a
separate acceptance test.
